%% kapa: exam=1 round=20 year=2009 period=1 form=A subject=civil_law subject_ko=민법 seq=40 number=40
%% source: Q-Net 감정평가사 기출문제 (cst00309) artlSeq=5024178 file=extracted/1008/5024178/2009년 제20회 감정평가사 1차시험 기출문제/1교시.hwp
%% conversion: kordoc 4.24.51 -> tools/gen1.py (auto); review=proofread-sonnet; flags=none
\begin{question}{40}
양도담보에 관한 설명으로 옳은 것은? (다툼이 있으면 판례에 의함)
\begin{choices}
  \item 양도담보권자는 사용ㆍ수익할 수 있는 정당한 권한이 있는 양도담보설정자에 대하여 자신이 사용ㆍ수익을 하지 못한 것을 이유로 임료 상당의 손해배상이나 부당이득반환의 청구를 할 수 있다.
  \item 피담보채권이 시효로 소멸하더라도 양도담보권은 소멸하지 않는다.
  \item 가등기담보 등에 관한 법률에 의하면 채무의 변제기가 지난 때부터 10년이 지난 이후에는 채무자는 채권담보의 목적으로 마친 소유권이전등기의 말소를 청구할 수 없다.
  \item 양도담보권은 특약이 없는 한 종물ㆍ과실에도 당연히 그 효력이 미친다.
  \item 채권의 담보목적으로 양도된 재산에 관한 담보권의 실행은 다른 약정이 없는 한 처분정산만이 허용된다.
\end{choices}
\end{question}
\answer{3}
